\writeSectionFrame{\presentationTitle}{Fazit}
\begin{frame}{Anwendungsfälle}
	\begin{itemize}
 		\item Kommunikation mit einfacher Schnittstelle
 		\item Unabhängige Weiterentwicklung von Subsystemen
 		\item Erhöhung der Sicherheit
 		\item Portierung
 		\item Isolierung von Anwendergruppen
 		\item Ausbauen von Funktionen
		\item Zusammenführung von Systemen
		\item Einteilung in Schichten
 	\end{itemize}
\end{frame}
 
\begin{frame}{Wann ist eine Fassade geeignet?}
	\begin{itemize}
 		\item Systeme hinter Fassade sind komplex
 		\item Fassade verbirgt viele Schnittstellen
 		\item Fassade bietet einen einfachen Zugriff
 		\item Schnittstellen ändern sich häufig
 		\item Subsysteme ändern sich häufig
 	\end{itemize}
\end{frame}

\begin{frame}{Weitere Aufgaben einer Fassade}
	\begin{itemize}
 		\item Validierungen
 		\item Authentifizierung
 		\item Logging
 		\item Steuerung von Transaktionen
 		\item Session-Verwaltung
 		\item Typ-Konvertierungen
 	\end{itemize}
\end{frame}

\begin{frame}{Besondere Fassadenarten}
	\begin{itemize}
 		\item Remote-Fassade
 		\item Message-Fassade
 	\end{itemize}
\end{frame}

\note{
	Die Remote-Fassade entkoppelt einen Client räumlich von den Subsystemen, d.h. die Subsysteme
	laufen auf entfernten Systemen und der Client soll davon nichts mitbekommen. Die Message-Fassade
	führt eine zeitliche Entkopplung durch, d.h. indem z.B. eine Message-Queue zwischen Client und
	Subsystem geschaltet wird. Diese Queue wird dann mit der Zeit abgearbeitet. 	
}

\begin{frame}{Remote-Fassade}
	\begin{center}
 		\includegraphics[scale=0.25]{\BILDERPFAD/fazit1.png}
 	\end{center}
\end{frame}

\note{
	Die Remote-Fassade soll einen Zugriff auf ein Netzwerk oder ein entferntes System so verstecken,
	dass der Client den Eindruck hat, er arbeite mit lokalen Objekten. Damit adressieren sie ein
	altes Problem der Kommunikation über Netzwerke: Service != Objekte. Objekte sind lokal und einfach
	zu erzeugen, während Services verteilte Systeme sind, deren Verwendung aufwändiger ist. Aber viele
	Frameworks erwecken den Eindruck, dass man nur mit lokalen Objekten arbeitet. Remotezugriffe müssen
	vor unberechtigten Zugriffen geschützt werden, sind unzuverlässiger als lokale Methodenaufrufe und 
	Vorkehrungen für Verbindungsabbrüche müssen geschaffen werden, verlangen eine Serialisierung und eine
	Deserialisierung und haben hohe Latenzzeiten.
}

\realworld{Web-Framework JavaServer Faces}{FacesContext-Klasse als Fassade}
\note{
    JSF ist ein Framework für die Entwicklung von Webanwendungen. Die FacesContext-Klasse agiert als Fassade für den Zugriff auf viele verschiedene JSF-Subsysteme, wie z.B. die Verwaltung von HTTP-Anfragen und -Antworten, Sitzungsmanagement und die Navigation. Entwickler können über FacesContext einfach auf diese Funktionen zugreifen, ohne sich mit den Details der darunterliegenden Implementierungen beschäftigen zu müssen.
}

\realworld{Spring Framework}{JdbcTemplate als Fassade}
\note{
    Spring bietet mit JdbcTemplate eine Fassade für die Arbeit mit JDBC (Java Database Connectivity). JdbcTemplate vereinfacht die Interaktion mit der Datenbank, indem es die Notwendigkeit für Boilerplate-Code wie das Öffnen und Schließen von Verbindungen, das Erstellen von Statements und das Verarbeiten von Ergebnismengen minimiert. Dadurch können Entwickler sich auf die eigentliche Logik konzentrieren, während JdbcTemplate die komplexen Details der Datenbankkommunikation handhabt.
}

\realworld{Hibernate}{Session- und EntityManager-Klassen als Fassaden}
\note{
    In Hibernate, einem ORM-Framework (Object-Relational Mapping), agieren die Session- und EntityManager-Klassen als Fassaden für die komplexe Interaktion zwischen der Anwendung und der Datenbank. Diese Klassen bieten Methoden zum Speichern, Laden, Aktualisieren und Löschen von Objekten, ohne dass sich Entwickler mit den SQL-Details auseinandersetzen müssen.
}

\realworld{SLF4J}{Fassade für Logging-Systeme}
\note{
    Bei SLF4J handelt es sich um ein Fassade, die verschiedene Logging-Systeme kaspelt, d.h. der Entwickler ist in seinem Client-Code unabhängig von dem tatsächlichen Logging-System, sondern er ruft nur die Methoden von SLF4J auf. 
}

\realworld{Apache Commons FileUtils}{Fassade für Dateisystemoperationen}
\note{
    Die FileUtils-Klasse aus Apache Commons IO bietet eine Fassade für Dateisystemoperationen. Sie stellt Methoden zum Kopieren, Verschieben, Löschen und Lesen von Dateien bereit, die den komplexen Umgang mit Dateisystemen und I/O-Streams erheblich vereinfachen. Entwickler können so Dateioperationen mit einfachen Methodenaufrufen durchführen, ohne sich um die Implementierungsdetails kümmern zu müssen.
}

\realworld{YouTube Java API}{Remote-Fassaden für Web-Zugriffe}
\note{
   Web-APIs wie die Java-API von YouTube sind Remote-Fassaden, die das wahrscheinlich sehr komplexe Subsystem vor dem Entwickler verstecken, einerseits weil er die komplexen Systeme nicht kennen muss, aber auch, weil er sie nicht kennen darf. 
}

\vorteil{Vereinfachung der Schnittstelle}
\note{
    Das Fassadenmuster bietet eine einfache und einheitliche Schnittstelle zu einem komplexen System, was die Bedienung erleichtert. Es ermöglicht es, mehrere komplexe Operationen durch einfache Methodenaufrufe auszuführen.
}

\vorteil{Reduzierte Abhängigkeiten}
\note{
    Durch die Verwendung einer Fassade interagieren Clients nur mit der Fassade und nicht direkt mit den komplexen Subsystemen. Dies reduziert die Kopplung und macht den Code modularer und leichter wartbar.
}

\vorteil{Verbesserte Wartbarkeit}
\note{
    Änderungen an den Subsystemen können vorgenommen werden, ohne den Client-Code zu beeinflussen, solange die Fassade nach außen hin gleich bleibt. Dies erleichtert die Wartung und Weiterentwicklung des Systems.
}

\vorteil{Zentralisierte Kontrolle}
\note{
    Eine Fassade kann die Interaktion mit verschiedenen Subsystemen zentralisieren und so eine kontrollierte und konsistente Verwendung dieser Subsysteme sicherstellen.
}

\vorteil{Erhöhte Lesbarkeit des Codes}
\note{
    Da der Code durch die Fassade vereinfacht wird, ist er leichter zu verstehen und zu lesen. Das macht es einfacher für neue Entwickler, sich in den Code einzuarbeiten.
}

\vorteil{Vielseitig einsetzbar}
\note{
    Die Fassade kann auf vielfältige Arten eingesetzt werden, z.B. Kapselung von Subsystemen, Zusammenführung fremder Systeme, Kapselung von Dritt-APIs, Kapselung von Fernzugriffen, Validierung usw. 
}

\vorteil{Realisierung verschiedener Zugriffsmöglichkeiten}
\note{
    Die Fassade ermöglicht auch die Implementierung verschiedene Zugriffsmöglichkeiten. So könnte man bei dem Beispiel mit dem Banken bspw. ein Abheben vom Girokonto und vom Sparbuch implementieren, die zwar unterschiedliche Vorgänge sind, aber für den Client transparent sind. 
}

\vorteil{Realisierung von Querschnittsfunktionalitäten}
\note{
    Querschnittsfunktionalitäten sind Funktionen, die mit der eigentlichen produktiven Logik nichts zu tun haben, z.B. Logging. Statt den Produktivcode mit Logging-Anweisungen aufzublähen, könnte dies in die Fassade ausgelagert werden. 
}

\vorteil{Lose Kopplung}
\note{
    Wenn gegen Interfaces programmiert wird, kann man auch innerhalb der Fassade eine extrem lose Kopplung erreichen und damit auch bspw. Subsysteme während der Laufzeit austauschen. 
}

\vorteil{Ermöglicht unabhängige Weiterentwicklungen}
\note{
    Die Untersysteme können dank der Kapselung hinter der Fassade unabhängig voneinander weiterentwickelt werden, weil sie sich untereinander auch nicht kennen. 
}

\nachteil{Zusätzliche Abstraktion}
\note{
    Das Einführen einer Fassade bedeutet eine zusätzliche Abstraktionsschicht, die den Code komplizierter machen kann, wenn sie nicht richtig eingesetzt wird. Dies kann zu einer unnötigen Komplexität führen.
}

\nachteil{Eingeschränkte Funktionalität}
\note{
    Eine Fassade kann dazu führen, dass nicht alle Funktionen der Subsysteme verfügbar sind, weil sie nur eine vereinfachte Schnittstelle bietet. Fortgeschrittene oder spezifische Funktionen der Subsysteme könnten durch die Fassade nicht abgedeckt werden.
}

\nachteil{Potenzielle Leistungsprobleme}
\note{
    Wenn die Fassade zu stark abstrahiert und viele Subsysteme gleichzeitig aufruft oder komplexe Logiken enthält, kann dies zu Leistungsproblemen führen. Jede zusätzliche Schicht kann potenziell die Performance beeinträchtigen.
}

\nachteil{Erhöhte Wartungskosten der Fassade}
\note{
    Die Fassade muss gepflegt und aktualisiert werden, wenn sich die darunterliegenden Subsysteme ändern. Dies kann zusätzlichen Wartungsaufwand bedeuten, insbesondere wenn die Fassade stark genutzt wird.
}

\nachteil{Versuchung, in Fassade zu viel Geschäftslogik zu implementieren}
\note{
    Die Geschäftslogik gehört natürlich in die Subsysteme und die Fassade hat einzig die Aufgabe, die Subsysteme zu koordinieren. Dennoch kann man dazu verführt werden zu viel Logik in die Fassadenklasse zu implementieren. 
}

\nachteil{Ändern sich Schnittstellen, muss Fassade geändert werden}
\note{
    Eine Fassade ergibt nur bei stabilen Schnittstellen Sinn. Unterliegen die Schnittstellen häufigen Änderungen, muss auch die Fassade angepasst werden, und diese Änderungen werden dann natürlich auf die Klienten übertragen. Und dann verliert die Fassade ihre Vorteile
}

\nachteil{Kompromisse notwendig, um alle Subsysteme unter einem Dach zu vereinen}
\note{
    Die Subsysteme sind oftmals sehr unterschiedlich. Um sie unter einem Dach zu vereinen, sind oftmals Kompromisse hinsichtlich der Funktionalitäten notwendig.
}

\nachteil{Man kann an der Fassade vorbei programmieren}
\note{
    Manchmal werden Entwickler dazu verführt, an einer Fassade vorbei zu programmieren. Auch damit verliert die Fassade sämtliche Vorteile. 
}

\nachteil{Fassade kann ein Gott-Objekt werden}
\note{
    Fassaden verführen dazu Gottklassen zu implementieren. 
}

\begin{frame}{An der Fassade vorbei programmieren (Gründe)}
	\begin{itemize}
 		\item Neue Clients können nicht zum Stichtag auf Fassade umgestellt werden
 		\item Entwickler kennen Fassade nicht
 		\item Gewohnheit oder Bequemlichkeit
 		\item Fassade ist technisch oder fachlich unzureichend
 	\end{itemize}
\end{frame}

\begin{frame}{An der Fassade vorbei programmieren (Lösungen)}
	\begin{itemize}
 		\item Einschränkung der Subsysteme (z.B. IP-Filter)
 		\item Sichtbarkeit der Subsysteme einschränken
 		\item Konstruktoren der Subsysteme verstecken
 		\item Durch Logging kann man feststellen, ob Klienten an Fassade vorbei auf Subsystem zugreifen
 	\end{itemize}
\end{frame}